% Segunda lectura — Bloque geométrico: transformaciones euclídeas y similaridades (enunciados originales de la guía)
\section{Bloque geométrico: transformaciones euclídeas y similaridades}\label{cap:08}\label{sec:bloque-euclideas}

{\small\vspace{1mm}\linkUnidad{03_reflexiones_similaridades}{Simulación interactiva: Unidad 5 — Reflexiones y similaridades}\par}


\section{Enunciados originales de la guía}

Los textos siguientes reproducen los enunciados de \texttt{Ejercicios.pdf}. Se normalizaron
únicamente los cortes de línea, la notación matemática y errores tipográficos evidentes; las
dos páginas originales se incluyen completas como anexo facsimilar al final del documento.

\begin{enunciado}[Definición 1]
La reflexión del punto \(a\) por un punto \(b\) es la operación que lleva el punto \(a\) al
punto que está sobre la recta \(\overline{ab}\), a la misma distancia de \(b\) que \(a\),
pero al otro lado de \(b\). Se llama \(\rho_b(a)\).
\end{enunciado}

\begin{enunciado}[Teorema 1]
\(\rho_a\rho_b\), la composición de dos reflexiones, es una traslación. Puede ser expresada
algebraicamente con la suma de un vector constante o geométricamente como el desplazamiento
según un segmento fijo.
\end{enunciado}

\begin{enunciado}[Teorema 2]
Si \(\rho_d=\rho_a\rho_b\rho_c\), entonces el polígono \(abcda\) es un paralelogramo.
\end{enunciado}

\begin{enunciado}[Observación]
Con esta operación puede definirse la suma y resta de vectores. ¿De qué modo? ¿Cómo se
relaciona esto con la propiedad de que sea inversiva la transformación, es decir, que al
aplicarla dos veces se anula?
\end{enunciado}

\begin{enunciado}[Definición 2]
La reflexión del punto \(a\) por la recta \(B\), \(\varphi_B(a)\), es la operación que lleva
el punto \(a\) al punto que está sobre la perpendicular a \(B\) que pasa por \(a\), a la misma
distancia que \(a\) de \(B\), pero al otro lado de la recta \(B\).
\end{enunciado}

\begin{enunciado}[Teorema 3]
Si \(\Theta_{AB}=\varphi_A\varphi_B\) y \(A\) es paralela a \(B\), entonces
\(\Theta_{AB}\) es una traslación. Puede ser expresada algebraicamente con la suma de un
vector constante o geométricamente como el desplazamiento según un segmento fijo.
\end{enunciado}

\begin{enunciado}[Teorema 4]
Si \(\Theta_{AB}=\varphi_A\varphi_B\) y \(A\) no es paralela a \(B\), entonces
\(\Theta_{AB}\) es una rotación.
\end{enunciado}

\begin{enunciado}[Teorema 5]
Si \(\Theta_{ABCD}=\Theta_{AB}\Theta_{CD}\), con \(\Theta_{AB}\) una traslación y
\(\Theta_{CD}\) una rotación, entonces \(\Theta_{ABCD}\) es una rotación. Además,
\(\Theta_{CDAB}=\Theta_{CD}\Theta_{AB}\) es también una rotación, pero es distinta.
\end{enunciado}

\begin{enunciado}[Definición 3]
Una homotecia con centro en el punto \(o\) y razón \(k\), nombrada \(\eta_o^k\), lleva un
punto \(a\) a otro sobre la recta \(\overline{oa}\), de forma que
\[
  \frac{d\bigl(o,\eta_o^k(a)\bigr)}{d(o,a)}=k,
\]
donde \(d(a,b)\) es la distancia entre el punto \(a\) y el punto \(b\). Los puntos \(a\) y
\(\eta_o^k(a)\) están al mismo lado del punto \(o\) en la recta \(\overline{oa}\).
\end{enunciado}

\begin{atencion}[Precisión sobre el texto original]
La razón entre distancias es siempre no negativa. Por eso, la redacción original corresponde
al caso \(k>0\). Para una razón firmada debe escribirse
\[
\eta_o^k(a)=o+k(a-o),
\qquad
\frac{d\bigl(o,\eta_o^k(a)\bigr)}{d(o,a)}=|k|.
\]
Punto e imagen están del mismo lado de \(o\) si \(k>0\), y en lados opuestos si \(k<0\).
\end{atencion}

\begin{enunciado}[Teorema 6]
Cualquier transformación de similaridad del plano puede expresarse como la composición
\(\varphi_A\varphi_B\eta_o^k\), donde el punto \(o\) es la intersección de las rectas
\(A\) y \(B\), y es el punto fijo de la transformación. Dado un triángulo y su imagen bajo
esta transformación puede hallarse el punto \(o\) con la construcción de Euler indicada en
la guía original.
\end{enunciado}

\begin{atencion}[Alcance del Teorema 6 original]
Tal como está escrito, el enunciado es demasiado general. La composición de dos reflexiones
y una homotecia describe una similaridad directa con punto fijo. Una similaridad inversa
requiere una cantidad impar de reflexiones, y una traslación no posee punto fijo finito.
\end{atencion}

\begin{enunciado}[Ejercicios]
Para cada una de las definiciones y teoremas expuestos, verificar la conclusión con un ejemplo
gráfico y construyendo un ejemplo algebraico con vectores. Pista para las reflexiones: para
hallar la proyección perpendicular de un punto sobre una recta puede usarse el producto
interno o construir la ecuación vectorial de la perpendicular usando el vector normal que
surge de la ecuación canónica de la recta.
\end{enunciado}

\section{Resolución completa de la guía}

La consigna de la guía no contiene una lista de problemas independientes. Solicita verificar
\textbf{cada definición y cada teorema} de dos maneras: mediante una construcción geométrica
y mediante un ejemplo algebraico con vectores. Las secciones siguientes constituyen esa
resolución punto por punto.

\begin{center}
\begin{tabularx}{\textwidth}{>{\bfseries}p{34mm} X X}
\toprule
Ítem de la guía & Verificación geométrica & Verificación algebraica\\
\midrule
Definición 1 & punto medio y segmentos iguales & \(\rho_b(x)=2b-x\)\\
Teorema 1 & desplazamiento constante & composición de dos reflexiones puntuales\\
Teorema 2 & diagonales que se bisecan & \(d=a-b+c\)\\
Definición 2 & perpendicular y distancias iguales & proyección mediante el vector normal\\
Teorema 3 & avance perpendicular de longitud \(2h\) & composición con normales paralelos\\
Teorema 4 & ángulo final igual a \(2\theta\) & producto de matrices de reflexión\\
Teorema 5 & mismo giro, distinto centro & resolución del punto fijo\\
Definición 3 & rayos desde el centro & \(\eta_o^k(x)=o+k(x-o)\)\\
Teorema 6 & escala y giro con centro común & \(S(x)=o+\lambda R(x-o)\)\\
\bottomrule
\end{tabularx}
\end{center}

\begin{idea}[Cómo reconocer que un ítem está resuelto]
Cada apartado presenta el dibujo o construcción, la propiedad geométrica que debe observarse,
la demostración vectorial y un ejemplo numérico. Los ejercicios del final son práctica
adicional y no sustituyen estas soluciones.
\end{idea}

